\section*{第 \ref{ch:g9:atoms-to-galaxies} 章——从原子到星系：宇宙的尺度}

\begin{solution}{exo:g9:atoms-to-galaxies:1}
\omterm{def:g9:atoms-to-galaxies:atom}{原子核}（$10^{-15}$ \unit{m}）、\omterm{def:g9:atoms-to-galaxies:atom}{原子}（$10^{-10}$）、\omterm{def:g7:states-of-matter:molecule}{分子}
（约 $10^{-9}$）、蚂蚁（$10^{-2}$，
若取较小的蚂蚁则为 $10^{-3}$）。
\end{solution}

\begin{solution}{exo:g9:atoms-to-galaxies:2}
极小而沉重、带正电的\omterm{def:g9:atoms-to-galaxies:atom}{原子核}位于中心；
几乎没有重量、带负电的\omterm{def:g9:atoms-to-galaxies:atom}{电子}云则远在其外。
大小方面：\omterm{def:g9:atoms-to-galaxies:atom}{原子}约 $10^{-10}$ \unit{m}，
\omterm{def:g9:atoms-to-galaxies:atom}{原子核}约 $10^{-15}$ ——相差十万倍的差距，
正是造就“体育场中一粒豌豆”这一比喻的原因。
\end{solution}

\begin{solution}{exo:g9:atoms-to-galaxies:3}
\omterm{def:g9:atoms-to-galaxies:atom}{电子}带负电，由 $-$ 极向 $+$ 极运动，
与约定方向相反。定律依然成立，
因为它们从不取决于谁在运动：电流、\omterm{def:g5:series-parallel-circuits:parallel}{节点}电流之和、\omterm{def:g8:ohms-law:resistance}{欧姆}
定律的图线与电费的计算，
读数完全一样，箭头依全球约定而定。
\end{solution}

\begin{solution}{exo:g9:atoms-to-galaxies:4}
取最接近的 10 的幂。身高：$10^{0}$ \unit{m}；
教室：$10^{1}$；地球直径：$10^{7}$。
\end{solution}

\begin{solution}{exo:g9:atoms-to-galaxies:5}
从\omterm{def:g9:atoms-to-galaxies:atom}{原子}到\omterm{def:g9:atoms-to-galaxies:atom}{原子核}：五个梯级。从你到地球：七个梯级。
从太阳到最近\omterm{def:g4:solar-system:star}{恒星}的距离（$10^{11}$）与（$4 \times 10^{16}$）之比：
相差五至六个梯级——约几十万倍。
\end{solution}

\begin{solution}{exo:g9:atoms-to-galaxies:6}
\omterm{def:g9:atoms-to-galaxies:atom}{原子}大部分是空的，直径约 $10^{-10}$ \unit{m}，
中心是 $10^{-15}$ 的\omterm{def:g9:atoms-to-galaxies:atom}{原子核}。把\omterm{def:g9:atoms-to-galaxies:atom}{原子}堆成桌子，空隙也随之而来。
\omterm{def:g3:solids-liquids-gases:solid}{固体}之所以坚实，是\omterm{def:g9:atoms-to-galaxies:atom}{电子}云的静电斥力
阻止它们相互重叠——空虚，却被牢牢守住。
\end{solution}

\begin{solution}{exo:g9:atoms-to-galaxies:7}
毫米，$10^{-3}$；\omterm{def:g9:atoms-to-galaxies:atom}{原子}，$10^{-10}$：相差七个梯级——
排成一行约有一千万个\omterm{def:g9:atoms-to-galaxies:atom}{原子}。
\end{solution}

\begin{solution}{exo:g9:atoms-to-galaxies:8}
\omterm{def:g8:speed-of-light:lightyear}{光年}：$10^{16}$ \unit{m}（约九点五千万亿
）。\omterm{ex:g5:stars-night-sky:milkyway}{银河}：$10^{21}$。已观测的最深天空：
接近 $10^{26}$。
\end{solution}

\begin{solution}{exo:g9:atoms-to-galaxies:9}
$10^{5}$ 粒豌豆：$0.01 \times 10^{5} = \qty{1000}{m}$ ——
宽一千米的“体育场”：比任何真实的都宏大，但
这幅图景依然成立——中心是一粒豌豆，
相邻的豌豆相距一千米。
\end{solution}

\begin{solution}{exo:g9:atoms-to-galaxies:10}
海洋中的水滴：约 $10^{21} \times 10^{5} = 10^{26}$
——因此，若照字面理解，旧说法便夸大了：海洋中的水滴
数比一滴水中的\omterm{def:g7:states-of-matter:molecule}{分子}数多出好几个梯级。
较诚实的说法——一滴水中的\omterm{def:g7:states-of-matter:molecule}{分子}数，多于一生
所能数过的任何东西——仍然成立；梯级让即便
心爱的说法也保持诚实。
\end{solution}

\begin{solution}{exo:g9:atoms-to-galaxies:11}
从\omterm{def:g9:atoms-to-galaxies:atom}{原子}往下约十级（自 $10^{0}$ 到 $10^{-10}$），
向上直到深空约二十六级，大致居中，
略偏向微小的一端。九年的学习足以让读者
诚实地触及所\omterm{def:g6:measurement-in-science:measuring}{测量}事物的
两端。
\end{solution}

\begin{solution}{exo:g9:atoms-to-galaxies:12}
答案因人而异——这道题本身就是那封信。
（仅为示范，举例如下：致读完“\omterm{def:g1:light-and-shadows:shadow}{影子}”一章的
一年级读者——“路面上那个黑色的孪生者，
终有一天会解释\omterm{def:g7:shadows-eclipses:solar}{日食}”；致读完“\omterm{def:g5:everyday-energy:energy}{能量}”一章的
五年级读者——“你追踪的那些\omterm{def:g5:everyday-energy:energy}{能量}链，终将
变成含\omterm{def:g9:kinetic-energy-safety:joule}{焦耳}的公式”；致读完\omterm{def:g8:ohms-law:resistance}{欧姆}定律的
八年级读者——“你画的图线，正是检验
定律的习惯，
这便是全部秘诀”。）
\end{solution}

\begin{solution}{pb:g9:atoms-to-galaxies:1}
\textbf{1.} 校园：第 $2$ 秒；地球：第 $7$ 秒；
\omterm{def:g2:moon-in-the-sky:moon}{月球}\omterm{def:g6:sun-earth-moon:axisorbit}{轨道}：第 $9$ 秒。
\textbf{2.} 第 $11$ 秒（距离 $1.5 \times 10^{11}$ \unit{m}）；
旁白：“毯子上的阳光，已是八
分钟前的旧光了。”
\textbf{3.} 大约第 $13$ 至 $16$ 秒：画面中
整个\omterm{def:g4:solar-system:family}{太阳系}只是一个光点，尚无\omterm{def:g4:solar-system:star}{恒星}——
说实话，接近漆黑的空旷，持续数秒：这是影片
最真实的镜头。
\textbf{4.} 已观测的最深天空，充满 $10^{26}$ \unit{m} 范围内的
星系——其最古老的光比地球本身还要古老。
\textbf{5.} 蚂蚁：向内第 $2$ 秒；细胞：第 $5$ 秒；
\omterm{def:g7:states-of-matter:molecule}{分子}：第 $9$ 秒；\omterm{def:g9:atoms-to-galaxies:atom}{原子}：第 $10$ 秒。
\textbf{6.} 穿过\omterm{def:g9:atoms-to-galaxies:atom}{原子}自身的空旷——
\omterm{def:g9:atoms-to-galaxies:atom}{电子}云与中心之间是五秒的虚无：“若这个\omterm{def:g9:atoms-to-galaxies:atom}{原子}
是一座体育场，我们已离开最高的看台，
仍在朝一粒豌豆飞去。”
\textbf{7.} \omterm{def:g9:atoms-to-galaxies:atom}{原子核}——坦率地说：“我们的梯子
到此为止，并非自然止步，而是因为这些
已是人类\omterm{def:g6:measurement-in-science:measuring}{测量}过的最小梯级；大厦
仍在延伸。”
\textbf{8.} 向外 $26$ 秒，向内 $15$ 秒：
一个四十一秒的宇宙。
\textbf{9.} 每一秒都使视野放大十倍——
相等的是\emph{比值}，而非相等的加数：梯子的
逻辑是：每一级都是一个因子，按因子计，
街道级的一步与星团级的一步一样大。
\textbf{10.} 例如：“无论看上去多满，几乎
这里的一切都是空的——物质与天空同样是
广阔虚无中稀疏的孤岛。”
\textbf{11.} 显微镜（\omterm{def:g8:lenses-images:families}{透镜}一章：两块\omterm{def:g8:lenses-images:families}{凸透镜}
用于观察微小之物）与望远镜（同一章，
收集遥远的微光）——两个方向的钥匙，
都由本书的光学原理磨制而成。
\textbf{12.} 例如：“九年，一把梯子：从
瓢虫可以数清的腿，到比脚下大地还古老的星系；
每一级都由任何人都可以核对的\omterm{def:g6:measurement-in-science:measuring}{测量}确定。
下一年，我们留下的那些问题——力与
运动、光的本质、\omterm{def:g9:atoms-to-galaxies:atom}{原子}的核心——开始
得到解答；带着这把梯子来吧。”
\end{solution}
